\section{Conclusion}

We classified Parkinson's disease from vertical ground reaction force
recordings with a Mamba state-space model on five subject-disjoint folds. On
both feet it reaches $75 \pm 4$\,\% accuracy, $0.75$ weighted F1 and $0.85$ AUC
at subject level, marginally ahead of the Random Forest, SVM-RBF, 1D-ResNet and
LSTM baselines run through the same evaluation code, though the Random Forest
remains better on the right foot alone. The result came from the input
representation rather than the architecture: the bidirectional extension of the
causal scan made no measurable difference at matched size.

The methodological finding may be the more useful one. Re-running the same
baselines with a window-wise split raises Random Forest accuracy from 73\,\% to
98\,\%, without touching the features, the models or the metrics. When one
recording yields more than a hundred correlated windows the split unit is not a
detail, and an accuracy quoted for this database cannot be interpreted until it
is stated. For the same reason we selected the final configuration on
validation AUC rather than test accuracy, which costs about three points but
makes the number an estimate of performance on new subjects rather than a
summary of our search.

More data would help more than more tuning. Every model we tried, classical or
deep, plateaued in the same narrow band, which points at the cohort rather than
the method: with 100 subjects one person is worth one accuracy point. Pooling
the dual-task recordings and other public VGRF cohorts is the obvious next step.
